\documentclass[12pt]{article}

% Core mathematics
\usepackage{amsmath,amssymb,amsthm}
\usepackage{mathtools}

% Diagrams
\usepackage{tikz-cd}
\usepackage{tikz}

% References
\usepackage{hyperref}
\usepackage[nameinlink,noabbrev]{cleveref}

% Graphics and tables
\usepackage{graphicx}
\usepackage{booktabs}
\usepackage{longtable}
\usepackage{array}

% Page layout
\usepackage[margin=1in]{geometry}
\usepackage{microtype}
\usepackage{enumitem}

% GrokRxiv sidebar
\usepackage{everypage}
\usepackage{xcolor}

\hypersetup{
  colorlinks=true,
  linkcolor=blue!55!black,
  citecolor=green!40!black,
  urlcolor=blue!60!black,
  pdftitle={Physical Realizability of Invertible Phase Classes},
  pdfauthor={Matthew Long and the YonedaAI Collaboration}
}

\newtheorem{theorem}{Theorem}[section]
\newtheorem{proposition}[theorem]{Proposition}
\newtheorem{lemma}[theorem]{Lemma}
\newtheorem{corollary}[theorem]{Corollary}
\theoremstyle{definition}
\newtheorem{definition}[theorem]{Definition}
\newtheorem{example}[theorem]{Example}
\newtheorem{counterexample}[theorem]{Counterexample}
\theoremstyle{remark}
\newtheorem{remark}[theorem]{Remark}

\newcommand{\Mic}{\mathfrak{Mic}}
\newcommand{\Inv}{\mathfrak{Inv}}
\newcommand{\FH}{\mathfrak{FH}}
\newcommand{\Cond}{\mathfrak{Cond}}
\newcommand{\IP}{\mathit{IP}}
\newcommand{\Gap}{\operatorname{Gap}}
\newcommand{\Spec}{\operatorname{Spec}}
\newcommand{\Th}{\operatorname{Th}}
\newcommand{\Pf}{\operatorname{Pf}}
\newcommand{\id}{\mathrm{id}}
\newcommand{\EST}{\textnormal{\textsc{Established}}}
\newcommand{\REAL}{\textnormal{\textsc{Realized}}}
\newcommand{\CAND}{\textnormal{\textsc{Candidate}}}
\newcommand{\PROP}{\textnormal{\textsc{Proposed}}}
\newcommand{\OPEN}{\textnormal{\textsc{Open}}}
\newcommand{\OBST}{\textnormal{\textsc{Obstructed}}}

\definecolor{grokgray}{RGB}{110,110,110}
\AddEverypageHook{%
  \ifnum\value{page}=1
    \begin{tikzpicture}[remember picture,overlay]
      \node[
        rotate=90,
        anchor=south,
        font=\Large\sffamily\bfseries\color{grokgray},
        inner sep=0pt
      ] at ([xshift=38pt,yshift=0.52\paperheight]current page.south west)
      {GrokRxiv:2026.07.physical-realizability\quad
       [\,math-ph\,]\quad
       14 Jul 2026};
    \end{tikzpicture}
  \fi
}

\title{Physical Realizability of Invertible Phase Classes:\\
Microscopic Witnesses, Bordism Targets, and Relative Charges}
\author{Matthew Long \\
\textit{The YonedaAI Collaboration} \\
\textit{YonedaAI Research Collective} \\
Chicago, IL \\
\texttt{matthew@yonedaai.com} $\cdot$ \url{https://yonedaai.com}}
\date{14 July 2026}

\begin{document}

\maketitle

\begin{abstract}
An abstract class in bordism, stable homotopy, or generalized cohomology is
not yet a quantum material.  A microscopic realization also needs a local
Hilbert space, an interaction, a symmetry action, a ground-state convention,
and a spectral gap bounded below independently of system size.  This paper
turns that observation into a realizability matrix.  Each row records an
explicit Hamiltonian, its symmetry, the theorem or exact calculation behind
its uniform gap, a microscopic invariant, and the status of its comparison
with an abstract class.

Three one-dimensional tests are worked out.  The class-D Kitaev chain has an
exactly flat bulk spectrum at its solvable point and realizes the nontrivial
fermionic $\mathbb Z_2$ phase.  A stack of BDI chains realizes the free
integer index, while the Fidkowski-Kitaev interaction identifies indices
differing by eight.  The spin-1 AKLT chain has a rigorous bulk gap and carries
the nontrivial projective $SO(3)$ edge representation.  These examples verify
specific low-dimensional rows.  They do not prove that every class counted
by a field-theory formula has a lattice representative.

We keep three spectra separate.  Freed and Hopkins classify deformation
classes of reflection-positive invertible field theories under stated
field-theoretic hypotheses.  Kubota constructs an operator-algebraic
$\Omega$-spectrum from invertible gapped quantum spin systems.  A condensed
phase spectrum remains a proposed organization of parameterized microscopic
families.  No comparison theorem presently identifies these objects.

Finally, for a family on $B$ with gapless locus $\Sigma$ and
$U=B\setminus\Sigma$, we state the exact hypotheses under which the
connecting morphism
$E^q(U)\to E^{q+1}(B,U)$ defines a relative class, localizes by excision, and
becomes a Thom class on $\Sigma$.  Its interpretation as a physical
transition charge still requires a specified microscopic comparison map.
The accompanying Haskell registry checks witness completeness, dimension
bookkeeping, and the degrees in this relative-charge construction.
\end{abstract}

\tableofcontents

\section{The realization problem}

Classifications of invertible phases appear in several mathematical forms.
Free fermions lead to operator K-groups.  Low-energy invertible field
theories lead to bordism and Anderson-dual groups.  Microscopic spin systems
can themselves be organized into homotopy types.  These outputs look similar
because stacking supplies addition and spatial suspension shifts degrees.
They are nevertheless built from different objects and different notions of
homotopy.

The question addressed here is deliberately concrete:

\begin{quote}
Given an abstract invertible class, is there a local symmetric Hamiltonian
whose thermodynamic phase maps to that class, with a gap controlled uniformly
in volume?
\end{quote}

The phrase ``maps to'' is essential.  A table in which two groups happen to
be isomorphic does not supply a map between them.  Nor does it prove that an
abstract generator has a representative satisfying a chosen microscopic
notion of locality.

\subsection{Spatial dimension and spacetime dimension}

Throughout the paper,

\begin{equation}
  d=\text{spatial dimension},
  \qquad
  n=d+1=\text{spacetime dimension}.
  \label{eq:dimension-rule}
\end{equation}

Hamiltonians and lattices are indexed by $d$.  The tangential symmetry type
in the Freed-Hopkins formula is indexed by $n$.  Thus every comparison row
must visibly apply the shift in \cref{eq:dimension-rule}.  In particular, a
one-dimensional chain has $d=1$ but contributes to a two-dimensional
spacetime theory with $n=2$.

\subsection{What this paper proves}

The main result is a verified portfolio, not a surjectivity theorem.

\begin{theorem}[Low-dimensional witness portfolio]
\label{thm:portfolio}
With the boundary conventions stated below, the following systems satisfy
the microscopic columns of the realizability matrix:

\begin{enumerate}[leftmargin=2.2em]
\item the solvable class-D Kitaev chain realizes the nontrivial free and
interacting one-dimensional fermionic $\mathbb Z_2$ class;
\item $r$ copies of the BDI chain realize free index $r$, and the
symmetry-preserving interacting classification depends on $r$ modulo eight;
\item the spin-1 AKLT chain realizes the nontrivial $SO(3)$-protected bosonic
one-dimensional phase detected by its spin-$1/2$ boundary representation.
\end{enumerate}

For each row there is an explicit finite-range interaction, an explicit
symmetry, a uniform bulk-gap calculation or theorem, and a microscopic
invariant.  The corresponding bordism or field-theory comparison is a
low-dimensional agreement using an effective-theory interpretation.  The
theorem does not construct a general map from microscopic phases to the
Freed-Hopkins spectrum.
\end{theorem}

The proof occupies \cref{sec:kitaev,sec:bdi,sec:aklt}.  The last sentence is
part of the theorem's scope, not a disclaimer added after the fact.

\subsection{Status language}

We use four statuses for realization questions.

\begin{definition}[Realization status]
Fix the microscopic rules, including locality, symmetry, stabilization,
boundary convention, and target invariant.

\begin{itemize}[leftmargin=2.1em]
\item \REAL{} means an explicit interaction, symmetry action, uniform-gap
verification, and invariant computation are all present.
\item \CAND{} means a model is explicit but at least one required witness is
missing.
\item \OPEN{} means no realization or obstruction theorem is asserted.
\item \OBST{} means a theorem rules out realization under the fixed rules.
\end{itemize}
\end{definition}

The absence of a known model gives status \OPEN{}, never \OBST{}.

\section{Three phase objects that must remain distinct}

\subsection{Microscopic stabilized phases}

Let $\Gamma$ be a lattice of bounded geometry.  Each site $x$ carries a
finite-dimensional Hilbert space $\mathcal H_x$, or a finite-dimensional
graded local algebra in a fermionic formulation.  A finite-range or
uniformly almost-local interaction $\Phi$ assigns a self-adjoint local
operator $\Phi(X)$ to each finite $X\Subset\Gamma$.

\begin{definition}[Microscopic phase datum]
A microscopic datum in spatial dimension $d$ is a tuple

\begin{equation}
  W=(\Gamma,\{\mathcal H_x\},\Phi,G,\alpha,\omega,\gamma,\nu),
  \label{eq:witness-tuple}
\end{equation}

where $G$ acts on the quasi-local algebra by $\alpha$, $\Phi$ is
$G$-invariant, $\omega$ specifies the ground-state sector, $\gamma>0$ is a
volume-independent bulk-gap lower bound, and $\nu$ is a phase invariant with
a declared codomain.  Fermionic parity is included in the symmetry data when
appropriate.
\end{definition}

For a periodic sequence of finite volumes $\Lambda_L$, the uniform-gap
condition used in the examples is

\begin{equation}
  \inf_{L\ge L_0}
  \left(E_{m_L+1}(H_{\Lambda_L})-E_{m_L}(H_{\Lambda_L})\right)
  \ge \gamma>0.
  \label{eq:uniform-gap}
\end{equation}

Here $m_L$ is the dimension of the prescribed low-energy sector.  For a
periodic chain with a unique ground state, $m_L=1$.  For an open AKLT chain,
the four edge states form the low-energy sector and the relevant bulk gap is
the gap above that sector.

Two data represent the same stabilized phase if, after adding atomic
product-state ancillas, they are joined by a symmetry-preserving path of
interactions that remains uniformly local and satisfies
\cref{eq:uniform-gap} along the path.  Invertibility means that stacking with
some second datum is gapped-homotopic, after stabilization, to the chosen
atomic phase.

\begin{remark}
An additive numerical invariant does not prove invertibility.  Invertibility
requires an inverse under stacking in the microscopic phase space.
\end{remark}

\subsection{Freed-Hopkins deformation classes}

Fix a spacetime symmetry type $(H_n,\rho_n)$, where
$\rho_n:H_n\to O_n$ has compact internal kernel.  Let $MTH$ be its
Madsen-Tillmann spectrum and let $I\mathbb Z(1)$ denote the indicated
Anderson-dual target.

\begin{theorem}[Freed-Hopkins, discrete topological sector]
\label{thm:fh}
\EST{}  Deformation classes of reflection-positive invertible extended
topological field theories of the fixed symmetry type are computed, in the
discrete topological sector, by the torsion subgroup

\begin{equation}
  \FH_n(H)
  :=[MTH,\Sigma^{n+1}I\mathbb Z(1)]_{\mathrm{tors}}.
  \label{eq:fh-group}
\end{equation}

The application of this result to lattice phases assumes the existence and
validity of an appropriate low-energy effective field theory.
\end{theorem}

The theorem is a field-theory classification.  It is not a theorem that each
element of \cref{eq:fh-group} has a finite-range lattice Hamiltonian.  The
broader extension to all nontopological reflection-positive invertible
theories is also not silently imported here.

\subsection{Kubota's microscopic Omega-spectrum}

\begin{theorem}[Kubota]
\label{thm:kubota}
\EST{}  There is an $\Omega$-spectrum $\IP_*$ constructed from invertible
gapped quantum spin systems on Euclidean spaces.  Its homotopy groups recover
the corresponding smooth homotopy groups of microscopic invertible systems.
The construction uses operator algebras, uniformly almost-local
interactions, relaxed lattices, stabilization by atomic systems, and sheaves
of smooth parameter families.  It also has variants for crystallographic
spatial symmetry.
\end{theorem}

The objects in this theorem include a chosen nondegenerate gapped GNS ground
state.  Compact Lie on-site symmetry yields the stated naive equivariant
variant.  This should not be upgraded without proof to a genuine
representation-graded equivariant spectrum.

\subsection{The proposed condensed spectrum}

A condensed moduli stack would evaluate parameterized Hamiltonian families
on profinite or more general condensed test objects.  Its invertible sector
could then be group-completed and stabilized.  This is a useful proposed
target for disorder and inverse-limit parameter spaces.  It is not Kubota's
smooth-manifold sheaf and it is not the bordism spectrum in
\cref{thm:fh}.

The current comparison situation is therefore

\begin{equation}
\begin{tikzcd}[column sep=large,row sep=large]
\Inv(\Mic_d^G)
  \arrow[r,dashed,"\text{partial microscopic packaging}"]
  \arrow[dr,dotted,swap,"\text{low-energy comparison}"]
& \IP_d^G
  \arrow[d,dotted,"\text{comparison open}"] \\
& \FH_{d+1}(H)
\end{tikzcd}
\qquad
\begin{tikzcd}[column sep=large]
\IP_d^G \arrow[r,dotted,"\text{no theorem}"]
& \Cond_d^G.
\end{tikzcd}
\label{eq:comparison-diagram}
\end{equation}

Solid arrows would require definitions on a common domain and proofs of
continuity, symmetry compatibility, stabilization compatibility, and
homotopy invariance.  The dotted arrows in \cref{eq:comparison-diagram} are
research questions.

\section{The realizability matrix}

\subsection{Witness columns}

\begin{definition}[Realizability row]
A realizability row for a proposed class $c$ contains the following fields:

\begin{enumerate}[leftmargin=2.2em]
\item spatial dimension $d$ and spacetime dimension $n=d+1$;
\item microscopic kinematics and boundary convention;
\item a local interaction $\Phi$;
\item an exact symmetry action and its square or extension law;
\item a uniform-gap witness, either an exact calculation or a cited theorem;
\item a microscopic invariant $\nu_{\mathrm{mic}}$;
\item an abstract target and a stated comparison status;
\item one of the four realization statuses.
\end{enumerate}
\end{definition}

The row is complete only if every required field is populated.  Numerical
evidence from finitely many lengths is not a uniform-gap theorem.

\subsection{The verified low-dimensional matrix}

\renewcommand{\arraystretch}{1.28}
\begin{longtable}{>{\raggedright\arraybackslash}p{0.14\textwidth}
                  >{\raggedright\arraybackslash}p{0.13\textwidth}
                  >{\raggedright\arraybackslash}p{0.17\textwidth}
                  >{\raggedright\arraybackslash}p{0.18\textwidth}
                  >{\raggedright\arraybackslash}p{0.21\textwidth}}
\toprule
System & Dimensions & Gap witness & Microscopic invariant & Abstract comparison \\
\midrule
\endfirsthead
\toprule
System & Dimensions & Gap witness & Microscopic invariant & Abstract comparison \\
\midrule
\endhead
Class-D Kitaev chain
& $d=1$, $n=2$
& Exact flat BdG spectrum at the solvable periodic point
& Pfaffian parity $-1$, equivalently one Majorana per open end
& Spin/Arf $\mathbb Z_2$ agreement in the effective-theory sector; no general spectrum map \\

BDI stack
& $d=1$, $n=2$
& Exact free bulk gap plus the Fidkowski-Kitaev symmetric interacting path for eight copies
& Free index $r$, interacting class $r\bmod 8$
& $\mathrm{Pin}^-$ and Arf-Brown-Kervaire $\mathbb Z_8$ agreement; comparison uses the field-theory interpretation \\

Spin-1 AKLT chain
& $d=1$, $n=2$
& Rigorous AKLT bulk-gap theorem
& Nontrivial projective $SO(3)$ edge representation
& Agreement with the $H^2(SO(3),U(1))\cong\mathbb Z_2$ sector; no general bordism realization theorem \\
\bottomrule
\end{longtable}

Each row has microscopic status \REAL{}.  Its last column states a more
limited comparison claim.  The word ``agreement'' does not mean that the
three spectra in \cref{eq:comparison-diagram} have been identified.

\begin{proposition}[Matrix verification criterion]
\label{prop:matrix-criterion}
A row has status \REAL{} if the interaction and symmetry are explicit, the
bulk gap satisfies \cref{eq:uniform-gap}, the microscopic invariant is
computed and stable under the chosen phase relation, and the row claims no
abstract comparison stronger than the supplied map or effective-theory
argument.
\end{proposition}

\begin{proof}
This is a contract rather than a classification theorem.  The first three
conditions produce a microscopic witness of the claimed phase.  The last
condition prevents a microscopic witness from being re-labelled as a
surjectivity theorem for a different classification object.
\end{proof}

\section{The class-D Kitaev chain}
\label{sec:kitaev}

\subsection{Hamiltonian and symmetry}

On a chain of $L$ complex fermions, let $c_j,c_j^*$ satisfy the canonical
anticommutation relations.  With periodic or open boundary convention as
specified, set

\begin{align}
H_D(\mu,t,\Delta)
={}&-\mu\sum_j\left(c_j^*c_j-\frac12\right)
-t\sum_j\left(c_j^*c_{j+1}+c_{j+1}^*c_j\right) \notag\\
&+\Delta\sum_j\left(c_jc_{j+1}+c_{j+1}^*c_j^*\right),
\label{eq:kitaev-hamiltonian}
\end{align}

where $t,\Delta\in\mathbb R$.  Fermion parity $(-1)^F$ is exact.  In BdG
form, particle-hole conjugation is intrinsic and squares to $+1$.  No
time-reversal symmetry is imposed, so the system is in class D.

Introduce Majorana operators

\begin{equation}
  a_j=c_j+c_j^*,
  \qquad
  b_j=\frac{c_j-c_j^*}{i}.
\end{equation}

At the solvable point $\mu=0$ and $\Delta=t>0$, an open chain has

\begin{equation}
  H_D=it\sum_{j=1}^{L-1}b_j a_{j+1},
  \label{eq:majorana-dimers}
\end{equation}

up to the conventional overall normalization.  The operators $a_1$ and
$b_L$ do not occur in \cref{eq:majorana-dimers}.  They form the two boundary
Majorana zero modes.

\subsection{Uniform periodic gap}

For a periodic chain, the BdG symbol may be chosen as

\begin{equation}
  h_D(k)=(-\mu-2t\cos k)\tau_z+2\Delta\sin k\,\tau_y.
  \label{eq:class-d-symbol}
\end{equation}

Here the Nambu spinor is
$\Psi_k=(c_k,c_{-k}^*)^{\mathsf T}$.  The matrices $\tau_y$ and $\tau_z$ are
Pauli matrices in this ordered particle-hole basis, with the pairing phase
chosen so that real $\Delta$ multiplies $\tau_y$.

Its positive quasiparticle energy is

\begin{equation}
  E(k)=\sqrt{(\mu+2t\cos k)^2+4\Delta^2\sin^2 k}.
  \label{eq:kitaev-dispersion}
\end{equation}

\begin{proposition}[Exact solvable-point gap]
\label{prop:kitaev-gap}
At $\mu=0$ and $\Delta=t>0$, every allowed periodic momentum satisfies
$E(k)=2t$.  Hence the periodic many-body bulk gap is bounded below by $2t$
for every $L$.
\end{proposition}

\begin{proof}
Substitution in \cref{eq:kitaev-dispersion} gives

\begin{equation}
  E(k)^2=4t^2\cos^2k+4t^2\sin^2k=4t^2.
\end{equation}

A quadratic BdG Hamiltonian diagonalizes into independent positive-energy
quasiparticles.  The least excitation energy is therefore $2t$, independent
of the momentum mesh and of $L$.
If one restricts the periodic Hilbert space to a fixed even-fermion-parity
sector, a single quasiparticle is excluded and the first allowed excitation
at this point has two quasiparticles and energy $4t$.  The stated $2t$
volume-independent lower bound is valid with or without that restriction.
\end{proof}

\subsection{The microscopic invariant}

At the particle-hole fixed momenta $k=0,\pi$, the class-D invariant can be
written as the sign of a product of Pfaffians.  In the convention of
\cref{eq:class-d-symbol},

\begin{equation}
  \nu_D=\operatorname{sgn}\bigl((\mu+2t)(\mu-2t)\bigr)\in\{+1,-1\}.
  \label{eq:pfaffian-invariant}
\end{equation}

Thus $|\mu|<2|t|$ gives $\nu_D=-1$.  The solvable point lies in this region.
The open-chain boundary Majoranas supply a second diagnostic of the same
phase.

\begin{corollary}[Class-D row]
The tuple consisting of \cref{eq:kitaev-hamiltonian}, fermion parity and BdG
particle-hole structure, periodic boundary conditions for the bulk gap, open
boundary conditions for the edge diagnostic, the lower bound $2t$, and
$\nu_D=-1$ is a \REAL{} microscopic row.
\end{corollary}

\subsection{Comparison status}

In spacetime dimension $n=2$, the corresponding invertible spin field theory
is the nontrivial Arf theory.  Both the microscopic interacting class-D
classification and this low-dimensional field-theory sector are
$\mathbb Z_2$.  The Kitaev chain is the standard microscopic representative
of the nontrivial element.

This agreement is strong evidence for the expected low-energy comparison.
It does not define the dotted arrow in \cref{eq:comparison-diagram} on all
microscopic phases, parameter families, or dimensions.

\section{BDI and the interaction reduction}
\label{sec:bdi}

\subsection{Free integer index}

Take $r$ copies of the real Kitaev chain.  Let the antiunitary symmetry $T$
act by complex conjugation in the real-fermion basis, so that

\begin{equation}
  T i T^{-1}=-i,
  \qquad
  T c_{j,\alpha}T^{-1}=c_{j,\alpha},
  \qquad
  T^2=+1.
  \label{eq:bdi-time-reversal}
\end{equation}

Together with BdG particle-hole structure this is class BDI.  The real
off-diagonal symbol has an integer winding number.  A stack of $r$ solvable
chains has

\begin{equation}
  \nu_{\mathrm{free}}=r\in\mathbb Z
\end{equation}

and carries $r$ Majorana zero modes at each open end.  Its periodic bulk gap
is again $2t$.

\subsection{Why the integer is not an interacting invariant}

\begin{theorem}[Fidkowski-Kitaev reduction]
\label{thm:fk}
\EST{}  For one-dimensional BDI fermions with the symmetry in
\cref{eq:bdi-time-reversal}, symmetry-preserving local interactions reduce
the free $\mathbb Z$ classification to $\mathbb Z_8$.  Eight boundary
Majorana modes can be gapped without a fermion bilinear and without breaking
the symmetry.  Correspondingly, eight copies of the nontrivial free chain
are connected to the trivial interacting phase by a symmetry-preserving
gapped path.
\end{theorem}

The interaction is local and quartic in the boundary Majoranas, and its bulk
extension supplies the gapped interpolation.  The cited construction, not a
finite-size numerical extrapolation, is the gap witness for the reduction.
For eight stacked chains, write $\gamma_a$ for the eight unpaired boundary
operators selected from the $a_j,b_j$ basis used in
\cref{eq:majorana-dimers}.  At an end where these $\gamma_a$ are $T$-even, a monomial
$\gamma_a\gamma_b\gamma_c\gamma_d$ with a real coefficient is also
$T$-even and preserves fermion parity.  The Fidkowski-Kitaev interaction is
built from such quartic terms, so the mechanism gaps the boundary multiplet
without violating \cref{eq:bdi-time-reversal}.

\begin{proposition}[Interacting BDI row]
\label{prop:bdi-row}
The microscopic invariant of the interacting BDI stack is

\begin{equation}
  \nu_{\mathrm{int}}=r\pmod 8.
  \label{eq:bdi-mod-eight}
\end{equation}

Rows $r$ and $r+8$ represent the same stabilized interacting phase, whereas
their free integer invariants differ.
\end{proposition}

\begin{proof}
The equivalence $r\sim r+8$ follows by stacking the Fidkowski-Kitaev gapped
trivialization of eight chains.  Their interacting classification also
distinguishes the remaining residue classes, for example through the
symmetry action and algebra of boundary degrees of freedom.  Hence the
quotient is precisely $\mathbb Z_8$.
\end{proof}

\begin{counterexample}[Free K-theory is not the universal interacting answer]
The free phases with indices $0$ and $8$ are distinct in the integer
classification.  By \cref{thm:fk}, they are equal after symmetry-preserving
interactions are admitted.  Therefore a free K-group cannot be identified
with the interacting microscopic phase group in general.
\end{counterexample}

\subsection{Pin structure and the abstract class}

The antiunitary relation $T^2=+1$ for the BDI fermion corresponds to the
$\mathrm{Pin}^-$ tangential structure in two spacetime dimensions.  The
Arf-Brown-Kervaire invariant has values in $\mathbb Z_8$, matching
\cref{eq:bdi-mod-eight}.  This is a particularly sharp low-dimensional
agreement between microscopic interaction reduction and bordism data.

The comparison still uses the low-energy field-theory interpretation in
\cref{thm:fh}.  It is not a construction of a natural equivalence between
Kubota's $\IP_*$ and the Freed-Hopkins mapping spectrum.

\section{The SO(3)-protected AKLT chain}
\label{sec:aklt}

\subsection{Interaction and symmetry}

At each site place the spin-1 irreducible representation of $SO(3)$.  Write
$\mathbf S_j=(S_j^x,S_j^y,S_j^z)$.  For neighboring sites let
$P^{(2)}_{j,j+1}$ be the orthogonal projection onto total spin two.  Since
$x=\mathbf S_j\cdot\mathbf S_{j+1}$ has eigenvalues $-2,-1,1$ in total-spin
sectors zero, one, and two,

\begin{equation}
  P^{(2)}_{j,j+1}
  =\frac16\left(
      (\mathbf S_j\cdot\mathbf S_{j+1})^2
      +3\mathbf S_j\cdot\mathbf S_{j+1}
      +2
    \right).
  \label{eq:aklt-projector}
\end{equation}

The AKLT Hamiltonian is

\begin{equation}
  H_{\mathrm{AKLT}}=\sum_j P^{(2)}_{j,j+1}.
  \label{eq:aklt-hamiltonian}
\end{equation}

Every term is a positive nearest-neighbor projection.  The diagonal on-site
$SO(3)$ action commutes with each term.

\subsection{Ground states, edges, and bulk gap}

The valence-bond construction represents each spin one as the symmetric
subspace of two virtual spin-$1/2$ variables and pairs neighboring virtual
spins into singlets.  On a periodic chain it gives the unique ground state.
On an open chain, one virtual spin-$1/2$ remains at each boundary, giving a
four-dimensional ground-state sector.

\begin{theorem}[AKLT bulk gap]
\label{thm:aklt-gap}
\EST{}  The one-dimensional AKLT interaction has a strictly positive bulk
spectral gap.  Equivalently, the periodic finite-volume gaps have a positive
lower bound for sufficiently large volume, and the open-chain spectrum has a
positive gap above its edge-state ground sector.
\end{theorem}

This is the imported gap theorem for the row.  Frustration freeness alone
would not have been enough.

\subsection{Projective boundary invariant}

The physical spin-1 representation descends from $SU(2)$ to a linear
representation of $SO(3)$.  A single virtual edge spin transforms in the
spin-$1/2$ representation of $SU(2)$.  The central element $-1\in SU(2)$
acts on it as $-\id$, so this action does not descend to a linear
representation of $SO(3)$.  It defines the nontrivial projective class

\begin{equation}
  [\omega_{\mathrm{edge}}]
  \in H^2(SO(3),U(1))\cong\mathbb Z_2.
  \label{eq:aklt-edge-class}
\end{equation}

\begin{proposition}[Stability of the AKLT edge class]
\label{prop:aklt-edge}
Along a uniformly gapped $SO(3)$-symmetric path of one-dimensional spin
interactions, the projective equivalence class of the boundary
representation is constant.
\end{proposition}

\begin{proof}[Proof sketch]
Quasi-adiabatic spectral flow transports the ground-state sector by a
quasi-local automorphism that intertwines the symmetry.  The induced
half-chain boundary representations are therefore unitarily equivalent up
to tensoring with linear on-site representations.  Such tensor factors do
not change the class in $H^2(SO(3),U(1))$.  The operator-algebraic version is
proved for compact symmetry groups by Bachmann and Nachtergaele.
\end{proof}

\begin{corollary}[AKLT row]
The interaction \cref{eq:aklt-hamiltonian}, diagonal spin-1 $SO(3)$ action,
AKLT bulk-gap theorem, periodic and open boundary conventions, and
nontrivial class \cref{eq:aklt-edge-class} form a \REAL{} microscopic row.
\end{corollary}

The same $\mathbb Z_2$ appears in the low-dimensional bosonic field-theory
sector with $SO(3)$ background.  As in the fermionic examples, this row
checks a generator and not the essential surjectivity of a general
comparison functor.

\section{Relative transition charge}
\label{sec:relative}

\subsection{The exact sequence of a pair}

Let $E$ be a multiplicative generalized cohomology theory represented by a
spectrum.  Let $B$ be a space in a category where the pair axiom and excision
hold, and let $\Sigma\subset B$ be closed.  Put
$U=B\setminus\Sigma$.  We assume the inclusions have been replaced by
cofibrant pairs when necessary, so that the cofiber model computes relative
cohomology.

The pair $(B,U)$ has the long exact sequence

\begin{equation}
\cdots\longrightarrow E^q(B)
\xrightarrow{j^*}E^q(U)
\xrightarrow{\partial}E^{q+1}(B,U)
\xrightarrow{i^*}E^{q+1}(B)
\longrightarrow\cdots.
\label{eq:long-exact-sequence}
\end{equation}

\begin{definition}[Relative class]
For $\nu\in E^q(U)$, define

\begin{equation}
  Q_\Sigma(\nu):=\partial\nu\in E^{q+1}(B,U).
  \label{eq:relative-charge}
\end{equation}

As algebraic topology this is an \EST{} connecting class.  Calling it a
physical transition charge is \PROP{} until $\nu$ is tied to a microscopic
family by a specified comparison map.
\end{definition}

\begin{proposition}[Extension criterion]
\label{prop:extension-criterion}
The class $Q_\Sigma(\nu)$ vanishes if $\nu$ is the restriction of a class in
$E^q(B)$.  Conversely, exactness implies that
$Q_\Sigma(\nu)=0$ precisely when $\nu$ lies in the image of $j^*$.
\end{proposition}

\begin{proof}
Exactness at $E^q(U)$ gives
$\ker\partial=\operatorname{im}j^*$.
\end{proof}

The converse concerns extension in the chosen cohomology theory.  It does
not say that a Hamiltonian family itself extends across $\Sigma$ while
remaining gapped.

\subsection{Excision near the gapless locus}

\begin{theorem}[Localization by excision]
\label{thm:excision}
Assume $B$ is paracompact Hausdorff and has the homotopy type required by
$E$.  Let $N$ be an open neighborhood of $\Sigma$ such that excision applies
to the pair $(B,U)$.  Then inclusion of pairs induces an isomorphism

\begin{equation}
  E^{q+1}(B,B\setminus\Sigma)
  \xrightarrow{\cong}
  E^{q+1}(N,N\setminus\Sigma).
  \label{eq:excision-isomorphism}
\end{equation}
\end{theorem}

\begin{proof}
Choose $N$ so that the closure of $B\setminus N$ is contained in the open
set $U$.  Excision removes $B\setminus N$ from both entries of the pair.
The resulting pair is $(N,N\setminus\Sigma)$.  Generalized cohomology
represented on cofibrant pairs satisfies this excision axiom.
\end{proof}

This theorem is the precise sense in which \cref{eq:relative-charge} is
supported near $\Sigma$.

\subsection{Thom identification}

Suppose now that $B$ is a smooth manifold and $\Sigma\hookrightarrow B$ is a
closed smooth submanifold of codimension $c$.  Let $\mathcal N\to\Sigma$ be
its normal bundle.  Assume $\mathcal N$ is $E$-oriented, meaning it has a
Thom class

\begin{equation}
  u_E\in\widetilde E^c(\Th(\mathcal N))
\end{equation}

whose restriction to every normal fiber is a generator in the required
sense.

\begin{theorem}[Relative class on the critical locus]
\label{thm:thom-charge}
Under the excision and orientation hypotheses above, a tubular neighborhood
and the Thom isomorphism give

\begin{align}
E^{q+1}(B,B\setminus\Sigma)
&\cong E^{q+1}(N,N\setminus\Sigma) \notag\\
&\cong \widetilde E^{q+1}(\Th(\mathcal N)) \notag\\
&\cong E^{q+1-c}(\Sigma).
\label{eq:thom-localization}
\end{align}

Thus $Q_\Sigma(\nu)$ determines a localized class of degree $q+1-c$ on
$\Sigma$.
\end{theorem}

\begin{proof}
The first isomorphism is \cref{thm:excision}.  Radial deformation in a
tubular neighborhood identifies the pair with the disk and sphere bundles
$(D\mathcal N,S\mathcal N)$, whose quotient is the Thom space.  Cup product
with $u_E$ is the Thom isomorphism and shifts degree by $c$.
\end{proof}

Without an $E$-orientation, the last line must use the corresponding twisted
cohomology.  Omitting that twist can change or destroy the purported charge.

\subsection{Linking spheres}

At a point $x\in\Sigma$ where the normal bundle is locally trivial, restrict
the Thom representative to a normal disk $D^c$ and its boundary $S^{c-1}$.
The resulting local class measures the obstruction on a linking sphere.
For ordinary cohomology this is the familiar local degree.  For K-theory or
another spectrum it is the corresponding generalized cohomology class.

\begin{remark}
A linking sphere must live in the full parameter space relevant to the
invariant.  For a band Hamiltonian this may include momentum together with a
tuning parameter.  A loop in the tuning parameter alone can miss the charge.
For example, an isolated codimension-three Weyl-type degeneracy is linked by
an $S^2$.  Its three normal coordinates may be three momentum components, or
two momenta together with a tuning mass in a parameterized two-dimensional
model.
\end{remark}

\subsection{Naturality required for physics}

Suppose a microscopic family $H$ on $U$ has an invariant
$\nu_{\mathrm{mic}}(H)$.  To interpret \cref{eq:relative-charge} physically,
one needs a natural transformation

\begin{equation}
  \kappa:\nu_{\mathrm{mic}}(H)\longmapsto \nu_E(H)\in E^q(U)
  \label{eq:comparison-map}
\end{equation}

that is invariant under the chosen gapped homotopies and compatible with
restriction to open sets.  Naturality then makes the connecting morphism
commute with the microscopic restriction maps.

\begin{proposition}[Conditional physical charge]
If the map \cref{eq:comparison-map} exists with these properties, then
$Q_\Sigma(\nu_E(H))$ is an invariant of the microscopic gapped family on
$U$ and is supported near the gapless locus by \cref{thm:excision}.
\end{proposition}

\begin{proof}
Homotopy invariance of $\kappa$ makes $\nu_E(H)$ a phase invariant.
Naturality of the long exact sequence transports this invariance through
$\partial$.  Excision supplies the support statement.
\end{proof}

The proposition is conditional.  A relative group by itself does not create
the microscopic map $\kappa$.

\section{Failures that a realizability claim must avoid}

\begin{counterexample}[Equal groups without a comparison]
Suppose two independently defined phase theories both produce
$\mathbb Z_8$.  Choosing generators gives many abstract group isomorphisms.
None of those choices proves that a microscopic Hamiltonian flows to a
particular field theory, respects families, or commutes with stacking.
\end{counterexample}

\begin{counterexample}[A finite-size gap table]
Assume exact diagonalization gives gaps above $0.2$ for lengths up to twenty.
The sequence could still decrease to zero at larger lengths.  The table is
evidence, not a witness for \cref{eq:uniform-gap}.
\end{counterexample}

\begin{counterexample}[Edge modes without a bulk convention]
An open finite chain can have nearly zero-energy boundary states for reasons
unrelated to a stable bulk phase.  Without a periodic bulk-gap statement and
a symmetry-stable boundary invariant, zero modes alone do not complete a
realizability row.
\end{counterexample}

\begin{counterexample}[Pointwise representatives without a family]
Choosing one Hamiltonian for each abstract class does not automatically give
a continuous map of phase spaces or spectra.  Parameter continuity and
uniform locality and gap bounds are additional data.
\end{counterexample}

\begin{remark}[Open is not obstructed]
If no lattice model is known for a bordism class, the correct status is
\OPEN{}.  Status \OBST{} requires a theorem such as an anomaly, locality, or
finite-dimensionality obstruction formulated for the same microscopic
rules.
\end{remark}

\section{What condensed organization contributes}

Condensed mathematics can organize families over profinite sets, disorder
hulls, and inverse limits.  It can impose descent on compatible local data
and make stacking functorial.  For the realizability problem, a condensed
test object could record how a portfolio of witnesses varies over a totally
disconnected compact parameter space.

It does not provide any missing witness in \cref{eq:witness-tuple}.  In
particular, sheaf descent does not prove the operator inequality behind a
gap, construct an inverse phase, or identify a low-energy field theory.

\begin{definition}[Proposed condensed realizability subfunctor]
For a condensed test object $S$, let $\mathcal R_d^G(S)$ consist, when
defined, of $S$-families of microscopic witnesses with a common locality
bound, a common positive gap, continuous local coefficients, and a locally
constant realizability status.  Pullback is by parameter restriction.
\end{definition}

This definition is \PROP{}.  A full construction must specify higher
isomorphisms, stabilization, and the topology on symmetry actions.  Even
after construction, a morphism
$\mathcal R_d^G\to\FH_{d+1}(H)$ remains a separate comparison problem.

\section{Executable witness registry}

The Haskell program under

\begin{center}
\path{src/physical-realizability/}
\end{center}

implements a finite registry rather than a symbolic classifier.  A row has
separate fields for $d$, $n$, interaction, symmetry, gap evidence,
microscopic invariant, abstract comparison, and status.  Validation checks
that $n=d+1$, that a \REAL{} row has every microscopic witness, and that no
row describes effective-theory agreement as a general spectrum equivalence.

The program also implements the degree bookkeeping

\begin{equation}
  q\xmapsto{\partial}q+1
  \xmapsto{\text{Thom, codim }c}q+1-c.
\end{equation}

It records whether extension forces the connecting class to vanish and
whether an $E$-orientation is available for the Thom step.

These are contract checks.  A Boolean value cannot prove the AKLT gap or the
Fidkowski-Kitaev interpolation.  The registry stores those results as cited
theorem witnesses and rejects rows in which the citation field is absent.

\section{Status ledger}

\begin{longtable}{>{\raggedright\arraybackslash}p{0.28\textwidth}
                  >{\raggedright\arraybackslash}p{0.15\textwidth}
                  >{\raggedright\arraybackslash}p{0.48\textwidth}}
\toprule
Statement & Status & Basis \\
\midrule
\endfirsthead
\toprule
Statement & Status & Basis \\
\midrule
\endhead
Freed-Hopkins classification of discrete reflection-positive invertible TFTs
& \EST{} & Stable homotopy and bordism theorem under its field-theory hypotheses \\

Kubota $\Omega$-spectrum for invertible gapped spin systems
& \EST{} & Operator-algebraic construction with smooth families and almost-local interactions \\

Kitaev class-D row
& \REAL{} & Explicit interaction, exact periodic gap, Pfaffian invariant, edge Majoranas \\

BDI mod-eight row
& \REAL{} & Explicit free stacks and Fidkowski-Kitaev interacting reduction \\

AKLT $SO(3)$ row
& \REAL{} & Explicit projector interaction, rigorous bulk gap, projective edge class \\

General microscopic to Freed-Hopkins comparison
& \OPEN{} & No map with all required functorial and analytic properties is supplied \\

Identification of Kubota and Freed-Hopkins spectra
& \OPEN{} & Domains and homotopies differ; no equivalence theorem is imported \\

Identification with a condensed phase spectrum
& \OPEN{} & Condensed spectrum remains a proposed construction \\

Connecting morphism in a pair
& \EST{} & Long exact sequence for generalized cohomology \\

Physical interpretation of every relative class
& \PROP{} & Requires a microscopic invariant and natural comparison map \\
\bottomrule
\end{longtable}

\section{Limitations and open problems}

\subsection{Restricted model portfolio}

All verified rows are one-dimensional.  They do not test chiral phases in
two spatial dimensions, intrinsic topological order, noninvertible phases,
or higher-form symmetry.  The AKLT row is a spin system, while the class-D
and BDI rows are fermionic.  Moving between these settings can involve a
choice of graded tensor product or Jordan-Wigner transformation and is not
treated as invisible.

\subsection{No general field-theory emergence theorem}

We do not prove that a gapped lattice model has a relativistic infrared
limit, that its low-energy theory is fully extended, or that this theory is
faithful to microscopic phase equivalence.  These are among the assumptions
needed to apply \cref{thm:fh} to lattice systems.

\subsection{No surjectivity}

Nothing here proves that every element of
$[MTH,\Sigma^{n+1}I\mathbb Z(1)]_{\mathrm{tors}}$ is realized by a local
Hamiltonian.  Nothing proves that every homotopy class in a proposed
condensed spectrum comes from Kubota's construction, or conversely.

\subsection{Boundaries and degeneracy}

Uniform bulk gaps and open-boundary spectral gaps are different statements.
We explicitly retain the low-energy edge sector in \cref{eq:uniform-gap}.
Boundary perturbations can split that sector while leaving the bulk gap
open.  A realizability row must state which notion it uses.

\subsection{Relative charge needs orientation and comparison}

The Thom identification in \cref{eq:thom-localization} can fail without an
$E$-orientation of the normal bundle.  Excision can fail if one works
outside a category of good pairs without a supported-cohomology
replacement.  Most importantly, a relative generalized cohomology class is
not automatically measurable until the microscopic invariant and response
map are specified.

\subsection{A useful next theorem}

A major next step would be a natural transformation from a controlled
microscopic invertible phase spectrum to an Anderson-dual bordism target,
defined on smooth parameter families and compatible with stacking,
suspension, spatial symmetry, and spectral flow.  Its value on the three
rows in this paper should recover the stated low-dimensional comparisons.
Constructing that map is \OPEN{}.

\section{Conclusion}

Physical realizability is a witness problem.  A class earns microscopic
status only after locality, symmetry, a thermodynamic gap, and an invariant
have been checked on an explicit system.  The Kitaev, BDI, and AKLT chains
pass that test in one spatial dimension.  Their agreement with spin,
$\mathrm{Pin}^-$, and $SO(3)$ field-theory data is informative and exact at
the level stated, but it is not a universal realization theorem.

The broader condensed program therefore has a clear boundary.  Condensed
objects can organize parameterized witnesses and their descent.  Kubota's
$\Omega$-spectrum supplies a genuine microscopic stable-homotopy object for
quantum spin systems.  Freed-Hopkins supplies a genuine bordism-based
classification of reflection-positive invertible field theories.  The maps
between these constructions remain mathematical work to be done.

\appendix

\section{Detailed witness sheets}

\subsection{Class D}

\begin{longtable}{>{\raggedright\arraybackslash}p{0.26\textwidth}
                  >{\raggedright\arraybackslash}p{0.65\textwidth}}
\toprule
Field & Witness \\
\midrule
Spatial and spacetime dimensions & $d=1$, $n=2$ \\
Local degrees of freedom & One complex fermion per site \\
Interaction & Nearest-neighbor hopping and pairing in \cref{eq:kitaev-hamiltonian} \\
Symmetry & Fermion parity; BdG particle-hole structure with square $+1$ \\
Bulk boundary convention & Periodic chain \\
Edge convention & Open chain \\
Gap & Exact $2t$ at $\mu=0$, $\Delta=t>0$ \\
Invariant & Pfaffian parity $\nu_D=-1$ \\
Edge diagnostic & One unpaired Majorana at each end \\
Microscopic status & \REAL{} \\
Abstract comparison & Nontrivial Arf/spin $\mathbb Z_2$ class under low-energy interpretation \\
Global comparison theorem & \OPEN{} \\
\bottomrule
\end{longtable}

\subsection{BDI}

\begin{longtable}{>{\raggedright\arraybackslash}p{0.26\textwidth}
                  >{\raggedright\arraybackslash}p{0.65\textwidth}}
\toprule
Field & Witness \\
\midrule
Spatial and spacetime dimensions & $d=1$, $n=2$ \\
Local degrees of freedom & $r$ real copies of the Kitaev chain \\
Interaction & Stacked quadratic chain plus allowed local quartic interactions \\
Symmetry & Antiunitary $T^2=+1$ and fermion parity \\
Free bulk gap & Exact $2t$ at the stacked solvable point \\
Interaction witness & Fidkowski-Kitaev gapped trivialization of eight copies \\
Free invariant & $r\in\mathbb Z$ \\
Interacting invariant & $r\bmod 8$ \\
Microscopic status & \REAL{} \\
Abstract comparison & $\mathrm{Pin}^-$ Arf-Brown-Kervaire $\mathbb Z_8$ agreement \\
Spectrum equivalence & \OPEN{} \\
\bottomrule
\end{longtable}

\subsection{AKLT}

\begin{longtable}{>{\raggedright\arraybackslash}p{0.26\textwidth}
                  >{\raggedright\arraybackslash}p{0.65\textwidth}}
\toprule
Field & Witness \\
\midrule
Spatial and spacetime dimensions & $d=1$, $n=2$ \\
Local degrees of freedom & Spin one per site \\
Interaction & Nearest-neighbor total-spin-two projectors \\
Symmetry & Diagonal linear $SO(3)$ action \\
Bulk convention & Unique periodic ground state \\
Edge convention & Four-dimensional open-chain ground sector \\
Gap & Rigorous AKLT bulk-gap theorem \\
Invariant & Nontrivial class in $H^2(SO(3),U(1))$ \\
Edge diagnostic & Projective spin-$1/2$ representation at each end \\
Microscopic status & \REAL{} \\
Abstract comparison & Nontrivial bosonic $SO(3)$ field-theory sector \\
General realization theorem & \OPEN{} \\
\bottomrule
\end{longtable}

\section{A direct check of the AKLT projector}

Let $x=\mathbf S_j\cdot\mathbf S_{j+1}$.  For two spin-one variables,

\begin{equation}
  x=\frac12\left(J(J+1)-2\cdot1(1+1)\right)
\end{equation}

in the total-spin-$J$ sector.  Hence $x=-2,-1,1$ for $J=0,1,2$.  The
polynomial

\begin{equation}
  p(x)=\frac{(x+2)(x+1)}6
\end{equation}

has values $0,0,1$ on those sectors.  Functional calculus therefore gives
$p(x)=P^{(2)}$, proving \cref{eq:aklt-projector}.  Positivity of every term
and annihilation of the valence-bond state follow immediately.  The
spectral gap remains a separate theorem because a sum of positive local
projectors need not be uniformly gapped.

\section{Relative-degree bookkeeping}

For reference, the exact sequence and Thom shifts are

\begin{equation}
\begin{tikzcd}[column sep=large]
E^q(B) \arrow[r,"j^*"]
& E^q(U) \arrow[r,"\partial"]
& E^{q+1}(B,U) \arrow[r,"i^*"]
& E^{q+1}(B)
\end{tikzcd}
\end{equation}

and, for an $E$-oriented normal bundle of rank $c$,

\begin{equation}
\begin{tikzcd}[column sep=large]
E^{q+1}(B,U)
  \arrow[r,"\text{excision}","\cong"']
& E^{q+1}(D\mathcal N,S\mathcal N)
  \arrow[r,"u_E^{-1}","\cong"']
& E^{q+1-c}(\Sigma).
\end{tikzcd}
\end{equation}

If $\nu=j^*\widetilde\nu$, then
$\partial\nu=\partial j^*\widetilde\nu=0$.  If the normal bundle is not
$E$-oriented, the final group must be replaced by the correctly twisted
group.  If $\Sigma$ is singular, one needs a stratified normal datum or a
supported-cohomology formulation rather than the smooth Thom statement used
here.

\section{Machine-checkable registry contract}

The finite registry validates the following implications:

\begin{align}
\texttt{REALIZED}
&\Longrightarrow
\texttt{interaction}\wedge\texttt{symmetry}, \\
&\phantom{\Longrightarrow{}}
\texttt{gap witness}\wedge\texttt{microscopic invariant}, \\
\texttt{dimension valid}
&\Longleftrightarrow n=d+1, \\
\texttt{extends globally}
&\Longrightarrow Q_\Sigma(\nu)=0, \\
\texttt{Thom degree}
&=q+1-c
\quad\text{when an orientation is present.}
\end{align}

It deliberately does not encode the false implication

\begin{equation}
  \texttt{same finite group}
  \Longrightarrow
  \texttt{equivalent phase spectra}.
\end{equation}

\section{Glossary of scope-sensitive terms}

\begin{longtable}{>{\raggedright\arraybackslash}p{0.24\textwidth}
                  >{\raggedright\arraybackslash}p{0.68\textwidth}}
\toprule
Term & Meaning in this paper \\
\midrule
Microscopic phase & Stabilized uniformly gapped homotopy class of local interactions with fixed symmetry rules \\
Invertible & Has a stacking inverse up to stabilized gapped homotopy \\
Uniform gap & Positive lower bound independent of volume, with a declared low-energy sector \\
Abstract class & Element of a stated K-theory, bordism, homotopy, or generalized cohomology group \\
Realized & Complete microscopic witness for the specific claimed invariant \\
Candidate & Explicit model with at least one missing witness \\
Open & Neither realization nor obstruction is asserted \\
Obstructed & Nonrealizability follows from a theorem under the same rules \\
Relative class & Image under the connecting morphism in a long exact sequence \\
Physical charge & Relative class equipped with a microscopic comparison and response interpretation \\
Condensed family & Proposed family organized on condensed test objects with uniform analytic bounds \\
\bottomrule
\end{longtable}

\begin{thebibliography}{99}

\bibitem{FreedHopkins}
D. S. Freed and M. J. Hopkins,
\emph{Reflection positivity and invertible topological phases},
Geometry and Topology \textbf{25} (2021), 1165 to 1330,
\url{https://arxiv.org/abs/1604.06527}.

\bibitem{Kubota}
Y. Kubota,
\emph{Stable homotopy theory of invertible gapped quantum spin systems I:
Kitaev's Omega-spectrum},
arXiv:2503.12618 (2025),
\url{https://arxiv.org/abs/2503.12618}.

\bibitem{KitaevWire}
A. Y. Kitaev,
\emph{Unpaired Majorana fermions in quantum wires},
Physics-Uspekhi \textbf{44} (2001), 131 to 136,
\url{https://arxiv.org/abs/cond-mat/0010440}.

\bibitem{KitaevPeriodic}
A. Kitaev,
\emph{Periodic table for topological insulators and superconductors},
AIP Conference Proceedings \textbf{1134} (2009), 22 to 30,
\url{https://arxiv.org/abs/0901.2686}.

\bibitem{FidkowskiKitaev2009}
L. Fidkowski and A. Kitaev,
\emph{The effects of interactions on the topological classification of free
fermion systems},
Physical Review B \textbf{81} (2010), 134509,
\url{https://arxiv.org/abs/0904.2197}.

\bibitem{FidkowskiKitaev2011}
L. Fidkowski and A. Kitaev,
\emph{Topological phases of fermions in one dimension},
Physical Review B \textbf{83} (2011), 075103,
\url{https://arxiv.org/abs/1008.4138}.

\bibitem{AKLT1987}
I. Affleck, T. Kennedy, E. H. Lieb, and H. Tasaki,
\emph{Rigorous results on valence-bond ground states in antiferromagnets},
Physical Review Letters \textbf{59} (1987), 799 to 802,
\url{https://doi.org/10.1103/PhysRevLett.59.799}.

\bibitem{AKLT1988}
I. Affleck, T. Kennedy, E. H. Lieb, and H. Tasaki,
\emph{Valence bond ground states in isotropic quantum antiferromagnets},
Communications in Mathematical Physics \textbf{115} (1988), 477 to 528,
\url{https://doi.org/10.1007/BF01218021}.

\bibitem{BachmannNachtergaele}
S. Bachmann and B. Nachtergaele,
\emph{On gapped phases with a continuous symmetry and boundary operators},
Journal of Statistical Physics \textbf{154} (2014), 91 to 112,
\url{https://arxiv.org/abs/1307.0716}.

\bibitem{ChenGuWen}
X. Chen, Z.-C. Gu, and X.-G. Wen,
\emph{Classification of gapped symmetric phases in one-dimensional spin
systems},
Physical Review B \textbf{83} (2011), 035107,
\url{https://arxiv.org/abs/1008.3745}.

\bibitem{PollmannTurnerBergOshikawa}
F. Pollmann, A. M. Turner, E. Berg, and M. Oshikawa,
\emph{Entanglement spectrum of a topological phase in one dimension},
Physical Review B \textbf{81} (2010), 064439,
\url{https://arxiv.org/abs/0910.1811}.

\bibitem{BMNS}
S. Bachmann, S. Michalakis, B. Nachtergaele, and R. Sims,
\emph{Automorphic equivalence within gapped phases of quantum lattice
systems},
Communications in Mathematical Physics \textbf{309} (2012), 835 to 871,
\url{https://arxiv.org/abs/1102.0842}.

\bibitem{Thiang}
G. C. Thiang,
\emph{On the K-theoretic classification of topological phases of matter},
Annales Henri Poincare \textbf{17} (2016), 757 to 794,
\url{https://arxiv.org/abs/1406.7366}.

\bibitem{TeoKane}
J. C. Y. Teo and C. L. Kane,
\emph{Topological defects and gapless modes in insulators and
superconductors},
Physical Review B \textbf{82} (2010), 115120,
\url{https://arxiv.org/abs/1006.0690}.

\bibitem{Brown}
E. H. Brown, Jr.,
\emph{Cohomology theories},
Annals of Mathematics \textbf{75} (1962), 467 to 484,
\url{https://doi.org/10.2307/1970209}.

\bibitem{AtiyahBottShapiro}
M. F. Atiyah, R. Bott, and A. Shapiro,
\emph{Clifford modules},
Topology \textbf{3} Supplement 1 (1964), 3 to 38,
\url{https://doi.org/10.1016/0040-9383(64)90003-5}.

\end{thebibliography}

\end{document}
